% Generated draft: requires mathematical and source review.
\documentclass[11pt]{article}
\usepackage[margin=25mm]{geometry}
\usepackage{fontspec}
\setmainfont{Latin Modern Roman}
\usepackage{amsmath,amssymb,mathtools,unicode-math}
\setmathfont{Latin Modern Math}
\usepackage{enumitem,needspace,xurl,hyperref}
\setlength{\parindent}{0pt}
\setlength{\parskip}{5pt}
\setlength{\emergencystretch}{3em}
\newcommand{\N}{\mathbb{N}}
\newcommand{\Z}{\mathbb{Z}}
\newcommand{\Q}{\mathbb{Q}}
\newcommand{\R}{\mathbb{R}}
\newcommand{\C}{\mathbb{C}}
\newcommand{\F}{\mathbb{F}}
\newcommand{\A}{\mathbb{A}}
\newcommand{\PP}{\mathbb{P}}
\newcommand{\E}{\mathbb{E}}
\newcommand{\eps}{\varepsilon}
\newcommand{\dd}{\,\mathrm{d}}
\begin{document}
\section*{UnsolvedMath ID: 20002908}
\textbf{Source problem code:} catalogue entry
\subsection{Definitions and mathematical statement}
The source asks for a precise relation between two gradings of products of Schur functions: the LLT grading and the Lapointe--Morse $t$-statistic. It also asks for the role of the affine Hecke algebra and the relation among $q$-Littlewood--Richardson coefficients, Macdonald polynomials, and $k$-Schur functions defined through generalized Kostka data. It supplies no explicit statistic, representation, parameter convention, or equality to be proved, so these are several connected research questions rather than one fully specified universal identity.

A coefficientwise formulation of a proposed grading reversal is as follows. In a fixed homogeneous degree let $F(X;q)=\sum_\lambda f_\lambda(q)s_\lambda(X)$ and $G(X;t)=\sum_\lambda g_\lambda(t)s_\lambda(X)$, with polynomial coefficients. An equality up to one global degree shift and parameter inversion means that there exists one integer $D$, independent of $\lambda$, such that
\[
F(X;q)=q^D G(X;q^{-1}),\qquad
[q^i]f_\lambda(q)=[t^{D-i}]g_\lambda(t)\quad\text{for all }\lambda,i.
\]
Zero coefficient polynomials must match. For each nonzero pair, necessary endpoint conditions are
\[
\min\deg f_\lambda+\max\deg g_\lambda=D,
\qquad \max\deg f_\lambda+\min\deg g_\lambda=D,
\]
with the same $D$ throughout; equality requires all reflected coefficients, not merely endpoints. Thus agreement after setting the grading parameter to one does not determine the proposed graded relation.

The aid reports exact agreement, after such a global inversion and shift, on Grojnowski--Haiman's nested rectangular locus: generalized Kostka coefficients for generalized Hall--Littlewood induction agree coefficientwise with the corresponding LLT $q$-Littlewood--Richardson coefficients. It also reports a $\operatorname{GL}_7$ $k$-split example of those authors disproving a universal extension. Neither the precise nested-rectangle hypotheses, the shift, nor the counterexample coefficient data are present here; these are qualified literature reports, not reproducible universal statements from this record alone.

The supplied normalization-consistent matrix relation can be specified independently of that proposed equality. Fix a degree and compatible finite index sets, and choose all transformed Macdonald and LLT conventions once, including any needed involution or parameter changes. Define matrices by
\[
s^{(k)}_\kappa=\sum_\lambda V^{(k)}_{\lambda\kappa}s_\lambda,
\qquad G_a=\sum_\lambda C_{\lambda a}s_\lambda,
\]
\[
H_\mu=\sum_\kappa A^{(k)}_{\kappa\mu}s^{(k)}_\kappa
=\sum_a W_{a\mu}G_a.
\]
Rows $\lambda$ index ordinary Schur functions, $\kappa$ indexes the admissible $k$-Schur basis, $a$ the LLT functions occurring in the expansion, and $\mu$ the admitted Macdonald labels. The last expansions presuppose membership in the relevant $k$-Schur span. In these conventions $C$ is the LLT-to-Schur $q$-LR matrix, $W$ the HHL LLT-weight matrix, $V^{(k)}$ the $k$-Schur-to-Schur matrix obtained recursively from the generalized-Kostka/$k$-split definition, and $A^{(k)}$ the Macdonald-to-$k$-Schur matrix. Equality of Schur coefficients gives
\[
V^{(k)}A^{(k)}=CW.
\]
This basis-factorization identity explains a coefficient connection once all data are defined; it neither proves universal equality of the two gradings nor supplies their combinatorial coefficients. The record gives no affine-Hecke-algebra module or action realizing the relationship, so that source subquestion remains unspecified and unresolved here.
\subsection{Short English statement}
The workshop asks how two graded Schur-product constructions and several related coefficient systems fit together. The supplied aid reports agreement on a restricted nested-rectangle family and a counterexample to universal agreement. It gives a precise coefficient-factorization identity and a test for a proposed global grading reversal, while leaving the detailed affine Hecke action and missing normalization data open to source review.
\subsection{Sources}
\begin{itemize}
\item Local registry: permanent catalogue entry, exact source code catalogue entry; \texttt{aim-special-functions-notes.json}, \texttt{notes[12]}.
\item AIM workshop \emph{Generalized Kostka polynomials}, Problem (13), \url{https://aimath.org/WWN/kostka/prob2.pdf}; original attribution to workshop suspicions and Morse--Lapointe's $t$-statistic. No proposer, year, remarks, or original literature citation is supplied.
\item Supplied attributions: Grojnowski--Haiman for the nested rectangular comparison and $\operatorname{GL}_7$ $k$-split counterexample; HHL for the LLT weights in a Macdonald expansion. No titles, theorem numbers, or complete authoritative publication references are included.
\item Supplied research attempt 1, registry updated 14 August 2026: partial result, low novelty confidence, machine-generated and not peer reviewed. It supplies the matrix relation and endpoint-reflection criterion. Verification identifies Problem 13 but ends before reproducing any further checked text.
\end{itemize}
\end{document}
